\documentclass[aspectratio=169,10pt,t]{beamer}
\usepackage{tema}
\usepackage{codice}
\usepackage{macros}
\begin{document}

\titolo{Ambiente e primo programma}

\begin{frame}{Obiettivi di oggi}
  \begin{itemize}
    \item Capire cosa succede dal file sorgente all'esecuzione
    \item Compilare ed eseguire un programma Java da terminale
    \item Riconoscere errori di compilazione ed errori di esecuzione
    \item Eseguire un programma da Visual Studio Code
    \item Vedere (solo un accenno) cos'è il debugger
  \end{itemize}
\end{frame}

\begin{frame}{Dal sorgente all'esecuzione}
  \begin{center}
    \texttt{Ciao.java} $\xrightarrow{\ \texttt{javac}\ }$ \texttt{Ciao.class} $\xrightarrow{\ \texttt{java}\ }$ risultato
  \end{center}
  \vspace{1em}
  \begin{itemize}
    \item \texttt{.java}: il testo che scriviamo (codice sorgente)
    \item \texttt{javac}: il compilatore, controlla e traduce in bytecode
    \item \texttt{.class}: il bytecode, eseguito dalla JVM con il comando \texttt{java}
  \end{itemize}
\end{frame}

\begin{frame}{Il primo programma}
  \codice{codice/Ciao.java}
  \vspace{0.5em}
  Il file si chiama \texttt{Ciao.java}: il nome del file coincide con quello della classe.
\end{frame}

\begin{frame}[fragile]{Compilare ed eseguire da terminale}
\begin{lstlisting}[language={},numbers=none]
$ javac Ciao.java
$ java Ciao
Ciao!
\end{lstlisting}
  \begin{itemize}
    \item \texttt{javac Ciao.java}: serve il nome del \textbf{file}, con \texttt{.java}
    \item \texttt{java Ciao}: serve il nome della \textbf{classe}, senza estensione
    \item Dopo la compilazione compare \texttt{Ciao.class} nella cartella
  \end{itemize}
\end{frame}

\begin{frame}[fragile]{Errore di compilazione}
  \codice{codice/ErroreCompilazione.java}
\begin{lstlisting}[language={},numbers=none]
ErroreCompilazione.java:3: error: ';' expected
    System.out.println("Ciao!")
                               ^
1 error
\end{lstlisting}
  Il compilatore indica la \textbf{riga} e il \textbf{tipo} di problema. Il programma non parte.
\end{frame}

\begin{frame}[fragile]{Errore di esecuzione}
  \codice{codice/ErroreEsecuzione.java}
\begin{lstlisting}[language={},numbers=none]
Exception in thread "main" java.lang.ArithmeticException: / by zero
	at ErroreEsecuzione.main(ErroreEsecuzione.java:5)
\end{lstlisting}
  Il programma compila, parte, e si ferma durante l'esecuzione.
\end{frame}

\begin{frame}{Due tipi di errore}
  \begin{center}
  \begin{tabular}{|l|l|l|}
    \hline
     & \textbf{Compilazione} & \textbf{Esecuzione} \\ \hline
    Quando & prima di partire (\texttt{javac}) & mentre gira (\texttt{java}) \\ \hline
    Il programma parte? & no & s\`i, poi si ferma \\ \hline
    Esempio & \texttt{;} mancante & divisione per zero \\ \hline
  \end{tabular}
  \end{center}
\end{frame}

\begin{frame}{Un accenno al debugger}
  Il debugger di VS Code esegue il programma \textbf{una riga alla volta} e mostra i valori delle variabili.
  \begin{itemize}
    \item \textbf{Breakpoint}: clic a sinistra del numero di riga (\texttt{F9}); il programma si ferma l\`i
    \item \textbf{Avvio}: pulsante \emph{Debug} sopra il \texttt{main}; \textbf{Step over} (\texttt{F10}): va avanti di una riga
    \item \textbf{Variables}: nel pannello \emph{Run and Debug}, i valori correnti delle variabili
  \end{itemize}
  \vspace{0.5em}
  Il programma \texttt{int x = 4; int y = 5; y = x + y;} visto passo passo:
  \begin{center}
  \begin{traccia}{2}
    dopo la riga & x & y \\ \hline
    1 & 4 & -- \\ \hline
    2 & 4 & 5 \\ \hline
    3 & 4 & 9 \\ \hline
  \end{traccia}
  \end{center}
  Dall'incontro 2 lo useremo davvero.
\end{frame}

\begin{frame}{Esercizi}
  \begin{esercizio}[Il tuo nome]
    Modifica \texttt{Ciao} perch\'e stampi il tuo nome al posto di \texttt{Ciao!}.
  \end{esercizio}
  \indizio{Cambia solo il testo tra virgolette.}
  \soluzione{\texttt{System.out.println("Mario");}}
\end{frame}

\begin{frame}[fragile]{Esercizi}
  \begin{esercizio}[Cosa stampa?]
    Prima di eseguire, scrivi cosa stampa \texttt{Righe}.
  \end{esercizio}
  \begin{columns}[T]
    \begin{column}{0.48\textwidth}
      \codice[numbers=none]{codice/Righe.java}
    \end{column}
    \begin{column}{0.48\textwidth}
      \indizio{\texttt{print} non va a capo, \texttt{println} s\`i.}
      \soluzione{\texttt{AB}, poi \texttt{CD}, poi \texttt{E}, ciascuna su una riga: la riga vuota \texttt{println()} chiude solo \texttt{CD}.}
    \end{column}
  \end{columns}
\end{frame}

\begin{frame}{Esercizi}
  \begin{esercizio}[Correggi l'errore]
    Copia \texttt{ErroreCompilazione.java}, compilalo e correggilo leggendo il messaggio.
  \end{esercizio}
  \indizio{Il messaggio dice \texttt{';' expected} alla riga 3.}
  \soluzione{Aggiungere il \texttt{;} alla fine di \texttt{System.out.println("Ciao!")}.}
\end{frame}

\begin{frame}{Esercizi}
  \begin{esercizio}[Errore di esecuzione]
    Esegui \texttt{ErroreEsecuzione}. In che riga si ferma? Perch\'e?
  \end{esercizio}
  \indizio{Guarda la riga indicata nel messaggio e i valori di \texttt{a} e \texttt{b}.}
  \soluzione{Riga 5: la divisione \texttt{a / b} con \texttt{b} uguale a 0 non \`e definita.}
\end{frame}

\begin{frame}[fragile]{Esercizio finale}
  \begin{esercizio}[Biglietto da visita]
    Scrivi un programma che stampi un biglietto da visita di 5 righe con una cornice di asterischi.
  \end{esercizio}
  \indizio{Una riga con \texttt{println} per ogni riga stampata.}
  \soluzione{\codice[numbers=none]{codice/Biglietto.java}}
\end{frame}

\begin{frame}{Errori frequenti}
  \begin{errore}
    \begin{itemize}
      \item Il nome del file non coincide con il nome della classe
      \item Maiuscole e minuscole: \texttt{System} non \`e \texttt{system}
      \item Virgolette o parentesi non chiuse
      \item \texttt{;} dimenticato a fine istruzione
      \item Eseguire con \texttt{java Ciao.java} o compilare con \texttt{javac Ciao}
    \end{itemize}
  \end{errore}
\end{frame}

\end{document}
